% Generated draft: requires mathematical and source review.
\documentclass[11pt]{article}
\usepackage[margin=25mm]{geometry}
\usepackage{fontspec}
\setmainfont{Latin Modern Roman}
\usepackage{amsmath,amssymb,mathtools,unicode-math}
\setmathfont{Latin Modern Math}
\usepackage{enumitem,needspace,xurl,hyperref}
\setlength{\parindent}{0pt}
\setlength{\parskip}{5pt}
\setlength{\emergencystretch}{3em}
\newcommand{\N}{\mathbb{N}}
\newcommand{\Z}{\mathbb{Z}}
\newcommand{\Q}{\mathbb{Q}}
\newcommand{\R}{\mathbb{R}}
\newcommand{\C}{\mathbb{C}}
\newcommand{\F}{\mathbb{F}}
\newcommand{\A}{\mathbb{A}}
\newcommand{\PP}{\mathbb{P}}
\newcommand{\E}{\mathbb{E}}
\newcommand{\eps}{\varepsilon}
\newcommand{\dd}{\,\mathrm{d}}
\begin{document}
\section*{UnsolvedMath ID: 30003748}
\textbf{Source problem code:} catalogue entry
\subsection{Definitions and mathematical statement}
\textbf{Source question.} For convolutional neural networks used in inverse problems, determine whether suitably discretized partial differential equation (PDE) motivated networks enforce stability, and determine how explicit prior knowledge and constraints should be incorporated.

\textbf{Possible mathematical framework, with unresolved specifications.} An inverse problem may be represented by a forward map \(F:X\to Y\), where \(X\) is a space of unknowns and \(Y\) is a space of observations. A reconstruction network is a map \(R_\theta:Y\to X\), with parameters \(\theta\). These symbols provide a framework for explaining the question; the source does not specify the spaces, the forward map, the network family, or the admissible parameters.

For example, if metrics \(d_X\) and \(d_Y\) are specified and stability means Lipschitz stability with respect to observation perturbations on an admissible set \(D\subseteq Y\), the corresponding property would be
\[
\exists C<\infty\quad\forall y,y'\in D,\qquad
 d_X\bigl(R_\theta(y),R_\theta(y')\bigr)\le C\,d_Y(y,y').
\]
This is only a possible interpretation of stability, not a condition selected or asserted by the source. In particular, the source does not specify whether such a constant must be uniform over network parameters, depth, discretization resolution, or other choices, or whether stability instead concerns numerical evolution, training, or another notion.

\textbf{Unresolved formulation.} The requested investigation is to identify PDE models, discretizations, architectural conditions, and ways of incorporating prior knowledge and constraints that yield the intended stability for reconstruction networks. No particular PDE, initial or boundary conditions, discretization rule, stability notion, class of priors, or constraint enforcement mechanism is supplied. Consequently, the record states a research question but does not determine a unique quantified mathematical conjecture. These missing specifications must be resolved before a precise theorem or conjecture can be extracted.
\subsection{Short English statement}
The question asks when convolutional reconstruction networks motivated by discretized PDEs are stable, and how prior knowledge and constraints should enter their construction. The supplied record leaves the meaning and scope of stability unspecified, so it does not determine a unique mathematical conjecture.
\subsection{Sources}
\begin{itemize}
\item Supplied local UnsolvedMath registry record, permanent ID \(30003748\), source problem code \texttt{catalogue entry}, title: \emph{Stable PDE-Informed Neural Networks for Inverse Problems}. Collection: Oberwolfach Reports Open Problems. No individual proposer is identified. The proposed year is recorded as 2018.
\item Supplied report citation: \emph{Mini-Workshop: Deep Learning and Inverse Problems (2019)}, Oberwolfach Reports, pp.\ 559--589, DOI 10.4171/owr/2018/11. Supplied source URL: \url{https://doi.org/10.4171/owr/2018/11}. The citation's 2019 designation differs from the recorded proposed year and the year appearing in the DOI; this discrepancy is unresolved.
\item Supplied literature reference: \emph{Oberwolfach Report 11}, \url{https://doi.org/10.4171/OWR/2018/11}.
\item Supplied literature reference: \emph{Stable architectures}, \url{https://arxiv.org/abs/1705.03341}.
\item Supplied literature reference: \emph{Convex and ODE stability}, \url{https://www.sciencedirect.com/science/article/pii/S0167278924001106}.
\item Literature assessment dated exactly \texttt{2026-08-21}: the supplied curation labels the problem open and reports that provably stable PDE-inspired architectures and conditional analyses exist, while the broad stability conjecture and its general mathematical formulation remain unresolved. This is supplied dated literature triage derived from the OWR curation, not an independently verified assessment. The record expressly requests independent verification of time-sensitive status and cited claims.
\item The registry labels the statement \texttt{corrected\_verified} and reports a manual semantic, typographic, and LaTeX review asserting that its clean statement is self-contained and preserves the mathematical target without unsupported additions. This is supplied review metadata, not independent mathematical verification of this draft or resolution of its formulation ambiguities.
\end{itemize}
\end{document}
